\documentclass[11pt]{article}
\usepackage[T1]{fontenc}
\usepackage{amsmath,amssymb,amsthm}
\usepackage[margin=1in]{geometry}
\usepackage{xurl}
\usepackage[hidelinks]{hyperref}
\newtheorem{theorem}{Theorem}
\newtheorem{lemma}[theorem]{Lemma}
\newtheorem{proposition}[theorem]{Proposition}
\newcommand{\adj}{\operatorname{adj}}
\newcommand{\tr}{\operatorname{tr}}
\title{Finite covers by monic one-variable standard algebras\\
\large A shorter local proof and a direct construction for one equation}
\author{}
\date{1 November 2026}
\begin{document}
\maketitle

The implication ``finitely presented, flat and formally unramified
implies formally \'etale'' is checked in \path{sag/etale-direct}.
This note studies the stronger conclusion: a finite principal cover by
monic one-variable standard \'etale algebras over the original base.
It removes the degree-lowering loop from the local proof, and gives a
direct mathematical construction for one equation, including
a bound on the number of charts. The shortened local argument is now
checked in \path{sag/etale-monic-local}; the finite-cover construction
remains mathematical. The final part gives the exact correspondence
with the formalisation and distinguishes the remaining obligations.

\part{Mathematics}

\section{The precise target}

All rings are commutative and may be the zero ring. A monic standard
\'etale algebra over $R$ has the form
\[
       (R[T]/(P))[1/(G(T)P'(T))],\qquad P\text{ monic}.
\]
The desired conclusion for an $R$-algebra $E$ consists of finitely many
$h_\alpha,c_\alpha\in E$, an identity
\[
                 \sum_\alpha c_\alpha h_\alpha=1,
\]
and monic standard presentations of each $E[1/h_\alpha]$ over $R$.
Both the identity and the actual algebra isomorphisms matter.

One cannot replace this by a single global generator. For example,
$\mathbf F_2^3$ is \'etale over $\mathbf F_2$, but any element has at
most two distinct coordinate values, and therefore cannot generate this
algebra. Nor does covering the closed fibre over a local base suffice:
for $R=k[t]_{(t)}$, the algebra $R\times R[1/t]$ has a whole component
invisible in its closed fibre.

\section{Surjective standard sources suffice}

The recent direct lifting theorem removes an entire issue from the
construction of the cover: once a standard source maps onto a proposed
chart, its kernel can be removed afterwards.

\begin{lemma}[Removing the kernel]\label{lem:kernel}
Suppose that $E$ is a finitely presented formally \'etale $R$-algebra,
and there is a surjection
\[
        u:S=(R[T]/(P))[1/(GP')]\longrightarrow E,
        \qquad P\text{ monic}.
\]
Then for some $F\in R[T]$,
\[
             E\simeq(R[T]/(P))[1/(GFP')].
\]
In particular the monic polynomial and its degree are unchanged.
\end{lemma}
\begin{proof}
Let $K=\ker u$. This ideal is finitely generated: choose preimages in
$S$ of the finitely many presentation generators of $E$; the defining
relations of $E$ and the equations recovering the generators of $S$
give a finite list of kernel generators.

Formal smoothness of $E$ lifts its identity to a section
$E\to S/K^2$. The composite $S\to E\to S/K^2$ and the quotient map
$S\to S/K^2$ agree modulo the square-zero ideal $K/K^2$. Formal
unramifiedness of $S$ makes these maps equal, whence $K=K^2$.
The determinant trick for a finitely generated idempotent ideal gives
$K=(e)$ with $e^2=e$. Consequently
\[
                    E=S/(e)\simeq S[1/(1-e)].
\]
Represent $1-e$ as $F(T)/(GP')^a$. Localising at it is the same as
inverting $F(T)$ in $S$.
\end{proof}

This lemma has no residue-field or flatness hypothesis. In the original
flat--unramified problem, its formal \'etaleness hypothesis is supplied
by the checked direct theorem.

\begin{proposition}[Reduction to square Jacobian algebras]
For the original flat--unramified problem it suffices to construct monic
standard covers of
\[
 S=\bigl(R[X_1,\ldots,X_r]/(F_1,\ldots,F_r)\bigr)[1/\Delta],
 \qquad \Delta=\det(\partial F_i/\partial X_j).
\]
\end{proposition}
\begin{proof}
The direct theorem constructs a surjection from such an $S$ to $E$,
and proves $E$ formally \'etale. The proof of
Lemma~\ref{lem:kernel}, through $K=(e)$, works for this $S$ too.
Thus $E\simeq S[1/(1-e)]$. Restrict a finite standard cover of $S$ to
this principal open. A further principal localisation of a standard
algebra is standard: express its denominator as a polynomial in the
standard generator divided by a power of the existing denominator.
The unit-ideal identity maps from $S$ to $E$.
\end{proof}

\section{Removing the degree-lowering loop from the local proof}

There is a substantial simplification even before attempting a direct
construction of the whole cover. Once a localisation is generated by
one integral element and a polynomial denominator, formal unramifiedness
supplies a simple monic relation without any calculation in the special
fibre.

\begin{lemma}[Monicisation of a simple relation]\label{lem:monicise}
Let $E$ be a finitely presented formally unramified $R$-algebra. Suppose
that $y\in E$ satisfies a monic polynomial $Q_0\in R[T]$, that $G(y)$
is invertible, and that the map
\[
                  R[T,1/G(T)]\longrightarrow E,
                  \qquad T\longmapsto y,
\]
is surjective. Then there is a monic polynomial $Q$ with
\[
                  Q(y)=0,\qquad Q'(y)\in E^\times.
\]
Consequently there is a surjection from a monic standard \'etale
algebra onto $E$. If $E$ is also formally smooth, it is itself monic
standard \'etale over $R$.
\end{lemma}
\begin{proof}
The kernel of the displayed map is finitely generated. Clearing the
denominators of finitely many generators gives polynomials
$F_1,\ldots,F_s\in R[T]$ such that
\[
             E=R[T,1/G]/(F_1,\ldots,F_s).
\]
We first show that the elements $F_j'(y)$ generate the unit ideal in
$E$. Put $C=E/(F_1'(y),\ldots,F_s'(y))$. There are two $R$-algebra
maps
\[
                  E\rightrightarrows C[\varepsilon]/(\varepsilon^2),
             \qquad y\longmapsto\bar y,\quad
                    y\longmapsto\bar y+\varepsilon.
\]
The second is well defined because
$F_j(\bar y+\varepsilon)=F_j(\bar y)+F_j'(\bar y)\varepsilon=0$,
and $G(\bar y+\varepsilon)$ is invertible, being invertible modulo a
square-zero ideal. The maps agree modulo $\varepsilon$. Formal
unramifiedness makes them equal, so $\varepsilon=0$, whence $1=0$ in
$C$. Effectiveness of the ideal quotient gives
\[
                         1=\sum_j b_j F_j'(y),\qquad b_j\in E.
\]
Express the finitely many $b_j$ as polynomial fractions in $y$ with
common denominator $G(y)^N$. There are $A_j\in R[T]$ with
\[
                G(y)^N=\sum_j A_j(y)F_j'(y).
\]
Set $H=\sum_j A_j F_j$. Then
\[
                      H(y)=0,\qquad H'(y)=G(y)^N.
\]
Choose a coefficient list for $H$ of length $k$, so that every recorded
term of $H$ has exponent less than $k$. Put
\[
                              Q=H+T^kQ_0^2.
\]
The two coefficient lists concatenate: the coefficients of $H$ come
first, followed by those of the monic polynomial $Q_0^2$. Thus $Q$ is
monic. Both $Q_0^2$ and its derivative vanish at $y$, so
\[
                  Q(y)=0,\qquad Q'(y)=H'(y)=G(y)^N,
\]
and the latter is a unit. This construction also covers $Q_0=1$ and
the zero algebra; no decision about a polynomial's actual degree is
needed. The alternative $Q_0^m+H$, with $m$ sufficiently large,
works when $Q_0$ has positive degree. Thus
\[
                 (R[T]/(Q))[1/(GQ')]\twoheadrightarrow E.
\]
If $E$ is formally smooth, apply Lemma~\ref{lem:kernel}.
\end{proof}

\begin{proposition}[An interface before localisation]\label{prop:interface}
Let $E$ be a finitely presented formally unramified $R$-algebra,
$y\in E$ integral over $R$, and $G\in R[T]$. Suppose that for every
$e\in E$ there merely exist $A\in R[T]$ and $n\geq0$ such that
\[
                         A(y)=G(y)^n e.
\]
If $E$ is formally smooth, or if $E$ is $R$-flat, then
$E[1/G(y)]$ is merely a monic standard \'etale $R$-algebra.
\end{proposition}
\begin{proof}
Principal localisation preserves finite presentation, formal
unramifiedness and formal smoothness. The displayed fraction
generation and the integral equation for $y$ give the hypotheses of
Lemma~\ref{lem:monicise} in the localisation. In the flat case the
direct flat--unramified theorem first supplies formal smoothness of
$E$. Every finite choice is made inside the truncation of the
conclusion. In particular $G(y)$ need not be invertible in $E$.
\end{proof}

\paragraph{Application to the existing local proof.}
Immediately after its determinant argument,
\path{zmt2-revised.tex} has
\[
                         E_g=R[y][1/G(y)],
\]
where $y$ belongs to a finite $R$-subalgebra and is therefore integral.
Lemma~\ref{lem:monicise} applies at exactly this point. It replaces the
construction of $\bar H$, the gcd/Nakayama degree-lowering loop, and
the special-fibre comparison used to find a suitable monic relation.
Its derivative is already invertible throughout $E_g$: no further
shrinking by a derivative is needed on the target. For a flat
unramified input, the checked direct theorem supplies formal smoothness,
and kernel removal finishes.

This does not remove the earlier construction of the finite subalgebra,
the residual monogenic factors, the residual idempotent, or the
determinant that produces $G$. It also does not by itself turn the
local charts into a finite cover over an arbitrary base. It does remove
one of the recursive parts of the local candidate construction.

\section{A cofactor lemma for recovering coordinates}

The following calculation does not require a free basis for the algebra.
The entries of $v$ below need only be a finite list of elements.

\begin{lemma}[Normalised eigenvector]\label{lem:eigen}
Let $M\in M_n(R)$ and let $y,v_0,\ldots,v_{n-1}$ belong to an
$R$-algebra $B$, with $v_0=1$ and $Mv=yv$. Set
\[
       P(T)=\det(TI-M),\qquad A(T)=\adj(TI-M),
       \qquad d_j(T)=A_{0j}(T).
\]
Then
\begin{align}
 P(y)&=0,\label{eq:root}\\
 A_{ij}(y)&=v_i d_j(y),\label{eq:cofactor}\\
 P'(y)&=\sum_{j=0}^{n-1}v_j d_j(y).\label{eq:trace}
\end{align}
Consequently the opens at $P'(y)d_j(y)$ cover $B[1/P'(y)]$.
On the $j$th such open,
\[
                      v_i=A_{ij}(y)/d_j(y).
\]
\end{lemma}
\begin{proof}
Write $L=yI-M$. The equation $Lv=0$, with $v_0=1$, expresses column
$0$ of $L$ as minus the sum of its other columns multiplied by the
corresponding $v_i$. Multilinearity and alternation of the determinant
give $\det L=0$. Apply the same column substitution to the cofactor
obtained by deleting row $j$ and column $i$. Every term with a repeated
column vanishes, and the remaining term, with its cofactor sign, is
$v_i$ times the cofactor obtained by deleting row $j$ and column $0$.
This gives~\eqref{eq:cofactor}, without cancelling any determinant or
assuming $P'(y)$ invertible.

Differentiating the determinant of $TI-M$ gives
$P'(T)=\tr A(T)$, by multilinearity in the columns. Substitute
\eqref{eq:cofactor} on the diagonal to obtain~\eqref{eq:trace}.
After inverting $P'(y)$ this is a unit-ideal certificate for the
$d_j(y)$.
\end{proof}

The needed matrix identities are therefore ordinary determinant
multilinearity and alternation. There is no rank test, basis extraction,
division by a nonunit, or appeal to prime ideals in this lemma.

\section{A direct finite cover for one equation}

\begin{theorem}\label{thm:basic}
Let
\[
       f(X)=\sum_{r=0}^n a_rX^r\in R[X],\qquad
       E=(R[X]/(f))[1/f'(x)].
\]
For $n\geq2$, there is a finite cover of $E$ by at most $2n^2$
monic standard \'etale algebras over $R$. Their monic polynomials have
degree $n$. The integer $n$ may be an upper bound for the degree:
no nonvanishing assumption on $a_n$ is required.
\end{theorem}

\subsection{Horner tails and finite matrices}

Put $D=f'(x)$, which is a unit in $E$, and define
\[
 t_i=a_i+a_{i+1}x+\cdots+a_nx^{n-i},\qquad 1\leq i\leq n.
\]
Let $V$ be the $R$-submodule spanned by $1,x,\ldots,x^{n-1}$.
Both $t_i$ and $xt_i$ preserve $V$. More precisely, for the exponents
needed here,
\begin{equation}\label{eq:horner}
 t_ix^j=
 \begin{cases}
  \displaystyle\sum_{r=i}^n a_rx^{r-i+j},&j<i,\\[2mm]
  \displaystyle-\sum_{r=0}^{i-1}a_rx^{r+j-i},&j\geq i.
 \end{cases}
\end{equation}
The second identity follows by subtracting $x^{j-i}f(x)=0$.
For $t_i$ use $0\leq j\leq n-1$, and for $xt_i$ use
$1\leq j\leq n$; every exponent on the right is at most $n-1$.

For each $i$ and $c\in\{0,1\}$ put
\[
                      y_{ic}=(x+c)t_i.
\]
Formula~\eqref{eq:horner} gives a concrete matrix $M_{ic}\in M_n(R)$
such that
\[
 M_{ic}(1,x,\ldots,x^{n-1})^{\mathsf t}
       =y_{ic}(1,x,\ldots,x^{n-1})^{\mathsf t}.
\]
Its row $j$ is obtained by adding the coefficient row for $t_ix^{j+1}$
to $c$ times the coefficient row for $t_ix^j$.
Define
\[
 P_{ic}(T)=\det(TI-M_{ic}),\qquad
 \delta_{ic}=P_{ic}'(y_{ic}).
\]
Each $P_{ic}$ is monic of degree $n$ and vanishes at $y_{ic}$ by
Lemma~\ref{lem:eigen}. The module $V$ need not be free, nor need it
equal $E$ as an $R$-module.

\subsection{The derivatives generate the unit ideal}

Here is the finite algebraic replacement for selecting the degree of
$f$ over a residue field. Work modulo $(a_{i+1},\ldots,a_n)$.
The first $i$ rows of $M_{ic}$ have support in the first $i$ columns.
The remaining diagonal block is triangular with diagonal
$-a_{i-1}$. The characteristic polynomial is consequently
\begin{equation}\label{eq:block}
 P_{ic}(T)=(T+a_{i-1})^{n-i}\Phi_{ic}(T),
\end{equation}
where the monic polynomial of the first block is
\[
 \Phi_{ic}(T)=(T-ca_i)^i+
       \sum_{r=0}^{i-1}a_ra_i^{i-1-r}(T-ca_i)^r.
\]
This formula has only nonnegative exponents, including when $i=1$.
It is the usual companion-matrix determinant, with scaled variable.
In this quotient $y_{ic}=a_i(x+c)$, and direct substitution gives
\[
 \Phi_{ic}(y_{ic})=a_i^{i-1}f(x)=0,
 \qquad a_i\Phi_{ic}'(y_{ic})=a_i^{i-1}D.
\]
Writing $w_c=a_i(x+c)+a_{i-1}$, differentiation of
\eqref{eq:block} therefore gives
\begin{equation}\label{eq:derivative-stratum}
 a_i\delta_{ic}=a_i^{i-1}w_c^{n-i}D
                \pmod{(a_{i+1},\ldots,a_n)E}.
\end{equation}
In particular $w_1-w_0=a_i$.

Let $J$ be the ideal generated by the $2n$ elements $\delta_{ic}$,
and use the constructive radical
\[
              \sqrt J=\{z:\text{ for some }N\geq1,\ z^N\in J\}.
\]
Its ideal laws follow from finite binomial expansions. We prove
$a_i\in\sqrt J$ by descending induction. Once the higher
coefficients belong to $\sqrt J$, equation~\eqref{eq:derivative-stratum}
and invertibility of $D$ give
\[
                 a_i^{i-1}w_c^{n-i}\in\sqrt J
                  \qquad(c=0,1).
\]
For $i=n$, this gives $a_n^{n-1}\in\sqrt J$, hence
$a_n\in\sqrt J$, because $n\geq2$.
For $1<i<n$, it gives $a_iw_0,a_iw_1\in\sqrt J$:
raise each product to a sufficiently large power and use radicality.
Subtracting gives $a_i^2\in\sqrt J$, hence $a_i\in\sqrt J$.
For $i=1$, it gives $w_0,w_1\in\sqrt J$, since $n-1\geq1$;
their difference is $a_1$.

Finally $D=\sum_{i=1}^n i a_i x^{i-1}$ belongs to $\sqrt J$.
Since $D$ is invertible, $1\in\sqrt J$, and hence $1\in J$.
This argument uses only finite ideal and power witnesses. It contains
no case distinction on whether a coefficient is zero or invertible.

\subsection{Cofactor charts and standard sources}

For a fixed pair $(i,c)$ abbreviate $y=y_{ic}$, $P=P_{ic}$,
$\delta=P'(y)$, and put
\[
 A(T)=\adj(TI-M_{ic}),\quad
 d_j(T)=A_{0j}(T),\quad p_j(T)=A_{1j}(T).
\]
The row with index $1$ exists because $n\geq2$.
Lemma~\ref{lem:eigen} supplies the identities
\[
      d_j(y)x=p_j(y),\qquad
      \delta=\sum_{j=0}^{n-1}x^j d_j(y).
\]
Set
\[
                    h_{icj}=\delta_{ic}d_j(y_{ic}).
\]
For each $(i,c)$, $\delta_{ic}^2$ belongs to the ideal generated by
its $h_{icj}$. Since the $\delta_{ic}$ generate the unit ideal, so
do all the $h_{icj}$. Explicitly, expand a sufficiently high power of
a unit-ideal identity for the $\delta_{ic}$ so that every summand
contains a square, and substitute the displayed cofactor identity.
There are $2n^2$ indices $(i,c,j)$; zero charts may be retained.

On $E[1/h_{icj}]$, both $\delta$ and $d_j(y)$ are invertible, and
$x=p_j(y)/d_j(y)$. Define the following polynomial, with no fractions:
\[
 H_j(T)=\sum_{r=1}^n r a_r p_j(T)^{r-1}d_j(T)^{n-r}.
\]
It satisfies
\[
                H_j(y)=d_j(y)^{n-1}D
\]
on this chart, so it is invertible there. Thus evaluation at $y$
gives a homomorphism from the monic standard algebra
\[
       S_{icj}=(R[T]/(P))[1/(P'd_jH_j)]
                  \longrightarrow E[1/h_{icj}].
\]
This homomorphism is surjective: it contains the generators
\[
 x=\frac{p_j(y)}{d_j(y)},\qquad
 D^{-1}=\frac{d_j(y)^{n-1}}{H_j(y)},\qquad
 h_{icj}^{-1}=\frac{1}{P'(y)d_j(y)}.
\]

The target is finitely presented and formally \'etale. Indeed a simple
root of $f$ lifts uniquely across a square-zero ideal by the correction
$a\mapsto a-f(a)/f'(a)$, and principal localisation preserves this
property. Lemma~\ref{lem:kernel} now turns the surjection into a
standard presentation, multiplying its denominator by one further
polynomial and keeping $P$ unchanged. This proves
Theorem~\ref{thm:basic}.

\subsection{Small bounds and a quadratic check}

For $n=0$, the inverted derivative is zero and $E$ is the zero ring.
For $n=1$,
\[
              E\simeq R[1/a_1],\qquad x=-a_0/a_1,
\]
which is standard with monic polynomial $T$ and denominator $a_1$.

For $f=aX^2+bX+c$, the choice $i=2$, $c=0$ in the construction gives
\[
 M=\begin{pmatrix}0&a\\-c&-b\end{pmatrix},\quad y=ax,
 \quad P(T)=T^2+bT+ac,
 \quad \adj(TI-M)=\begin{pmatrix}T+b&a\\-c&T\end{pmatrix}.
\]
Here $P'(y)=2ax+b=D$ is already a unit. The two cofactor denominators
are $ax+b$ and $a$, and
\[
                    D=(ax+b)+ax.
\]
They give exactly the two opens in the direct quadratic example in
\path{zmt2-revised.tex}. Thus the quadratic construction extends to
arbitrary degree through finite matrices and a bounded induction.

\section{What this does and does not settle}

The construction proves the monic-cover theorem directly for a single
equation with invertible derivative. It works over the original ring,
including nonreduced rings and finite fields. It uses neither ZMT,
residue-field factorisation, selected primes, nor a dynamical
interpretation of the local proof.

For a general square Jacobian algebra the unresolved task is still to
obtain such one-variable charts. Gaussian elimination on the Jacobian
does not by itself eliminate polynomial variables. Similarly, replacing
a system by a tower of one-variable extensions would still require
an argument turning that tower into a finite monic standard cover.
Monogenicity is not preserved by arbitrary compositions: over
$\mathbf F_2$, the tower $\mathbf F_2\to\mathbf F_2^2\to\mathbf F_2^4$
already illustrates the problem.

Lemma~\ref{lem:eigen} suggests a concrete algebraic interface for a
possible extension. For each of finitely many candidates, provide:
\begin{enumerate}
\item a finite vector $v=(1,x_1,\ldots,x_r,\ldots)$ of elements of the
target algebra;
\item a matrix $M$ over $R$ and an element $y$ with $Mv=yv$;
\item a unit-ideal identity for the corresponding
$\det(TI-M)'|_{T=y}$, over all the candidates.
\end{enumerate}
The cofactor construction then expresses every chosen algebra generator
as a rational function of $y$ on finitely many opens. It provides
surjective standard sources, and Lemma~\ref{lem:kernel} finishes.
Any localisation generators must also be included or have their
inverses explicitly supplied, as $D^{-1}$ was above.

This interface is sufficient, not an already proved existence theorem.
In particular a finite spanning module with arbitrary multiplication
matrices can give a characteristic polynomial with repeated roots.
Formal unramifiedness of the target does not make the derivative of
such an arbitrary characteristic polynomial invertible. The Horner
construction above succeeds because the block calculation controls
those extra factors and the coefficient induction covers their failures.
An analogue of that calculation for several variables is the main
research question.

\part{Connection with the checked development}

\section{Already available}

\begin{itemize}
\item \path{sag/etale-direct/FlatUnramifiedEtale.seg} proves the
arbitrary-base lifting theorem from finite presentation, flatness and
formal unramifiedness.
\item \path{sag/etale-direct/UnramifiedEtaleSource.seg} constructs the
surjective square Jacobian source. \path{JacobianLifting.seg} supplies
the square-zero Newton correction.
\item \path{sag/etale-local/EtaleCase.seg}, together with
\path{IdempotentIdeal.seg}, proves the kernel-removal mechanism in
Lemma~\ref{lem:kernel}. Its application needs finite kernel generators,
fraction representatives and the explicit source presentation. These
bridges are now constructed for the local shortcut in
\path{sag/etale-monic-local}.
\item \path{sag/etale-dynamical-global/StandardChartLocalisation.seg}
proves closure of standard presentations under further principal
localisation, with the algebra universal property.
\item The determinant, adjugate, characteristic-polynomial, polynomial
and localisation developments provide the underlying operations.
In particular \path{Adjugate.seg} defines cofactors by column
replacement and proves their linearity and adjugate identities.
\end{itemize}

These are reusable ingredients. They do not mean that the new finite
cover or its unit-ideal certificate has already been checked.

\section{The checked shortened local proof}

The development \path{sag/etale-monic-local} proves
Lemma~\ref{lem:monicise} and Proposition~\ref{prop:interface}.
Its public endpoint is
\begin{quote}\small
\path{ShortLocalEtale.Over[R,setR,cs].At[E,y,full,unique]}\\
\path{theorem fp integral smooth}, or\\
\path{flat-theorem fp integral flat}.
\end{quote}
Here $G$ is the polynomial represented by \path{cs}; \path{full} is
exactly the fraction-generation hypothesis of
Proposition~\ref{prop:interface}; and \path{unique} is formal
unramifiedness. Both \path{fp} and \path{integral} are truncated:
neither presentation generators nor the monic integral equation have
to be chosen outside the conclusion.

The output merely consists of a localisation witness for $E[1/G(y)]$,
including its algebra universal property, a coefficient list specifying
a monic $Q$, a polynomial denominator $D$, and a standard presentation
\[
                 E[1/G(y)]\simeq(R[T]/(Q))[1/(DQ')].
\]
Standard presentations are expressed by the checked
\path{StandardEtale.Points.Standard} universal property, not just
by a surjective map or an equality of sets of points.

The construction has the following correspondence with the proof.
\begin{itemize}
\item \path{FiniteSurjectionKernel.seg} constructs finite kernel
generators for a surjection between finitely presented algebras.
It is used both for $R[T,1/G]\to E$ and for the final standard source.
\item \path{PolynomialLocalisation.seg} and
\path{PolynomialFractions.seg} construct the source presentation,
maps and fraction arithmetic.
\item \path{MonicDualTest.seg} and
\path{MonogenicDerivativeIdeal.seg} construct the dual-number map
over the quotient by the derivatives and obtain
$1=\sum_j b_jF_j'(y)$. These witnesses are proved from unramifiedness,
not added as hypotheses. This step uses the direct dual-number test
rather than an application of \path{DiagonalJacobian.seg}.
\item \path{SimplePolynomialRelation.seg} clears the finite family
of fraction coefficients and proves $H(y)=0$ and $H'(y)=G(y)^N$.
\path{MonicRelation.seg} then constructs exactly $H+T^kQ_0^2$ by
coefficient-list concatenation and proves its value and derivative.
\item \path{MonogenicStandard.seg} builds the monic standard
surjection, constructs its finite kernel and calls the unchanged
\path{EtaleCase.seg} kernel-removal theorem.
\item \path{PrincipalEtale.seg},
\path{LocalMonogenicStandard.seg} and \path{ShortLocalEtale.seg}
provide localisation and the public interface above. The flat variant
calls the unchanged \path{FlatUnramifiedEtale.seg} theorem.
\end{itemize}

No field, prime, ZMT, classical-logic or choice principle is a new
hypothesis. The ambient univalence, replacement and surjection
parameters are the same as in the reused development. The old local
proof file is not rewritten: the new theorem supplies its replacement
ending after the determinant and fraction-generation data have been
constructed. The earlier construction of those data remains necessary.

A fresh run checks 6,795 declarations with zero failures in 2.78 seconds;
679 declarations in 12 proof modules are new. The opening audit covers
all 657 named new declarations; each opens only its declared list,
with at most 33 definitions opened and at most 0.01 seconds per new
declaration in this run. The module check passes all 168 modules.

Regressions include repeated initial monic relations, redundant and
empty presentations, arbitrary already-standard algebras, and
localisations of $\mathbf Z/4\mathbf Z$ at $1$, $2$ and $0$.
A concrete two-denominator calculation checks the cleared exponent
$1+2=3$. Mutations removing the square from $Q_0^2$, changing the
tangent coefficient from $1$ to $0$, and shifting the cleared exponent
are all rejected by the checker. These mutation checks use temporary
copies; the development itself remains unchanged.

\section{The separate finite-cover construction}

The one-equation cover theorem has not been formalised in this work.
It remains a separate, bounded target:
\begin{enumerate}
\item Prove the normalised-eigenvector cofactor identities and the
trace formula for the derivative of the characteristic polynomial.
Use the existing column conventions carefully: the mathematical
matrices here are written with row index first, whereas Unit's
determinant API stores columns first.
\item Construct the Horner multiplication matrices and their
eigenvector equations. Prove the block determinant formula
\eqref{eq:block}, including $i=1$, zero leading coefficients and
the empty lower block.
\item Convert the descending radical argument into finite power and
ideal-combination certificates. This needs no prime-existence axiom
and no choice of an actual degree for $f$.
\item Construct the $2n^2$ sources, prove surjectivity from the
cofactor expressions, instantiate kernel removal, and export the
finite cover with the original monic characteristic polynomials.
\end{enumerate}

The new work is mainly the two determinant calculations and their
polynomial bookkeeping. The previous lifting and idempotent-kernel
proofs can be reused. A successful formalisation of this milestone
would not yet complete the several-variable cover theorem.

As an informal arithmetic check, the matrix construction was tested
on 2,136 characteristic-polynomial block identities and 806 instances
with a root and an invertible derivative. The rings were
$\mathbf Z/m\mathbf Z$ for $m\in\{2,3,4,5,6,8\}$. The cofactor identities and
the unit ideal generated by the proposed chart denominators passed.
These finite tests support the calculations; they do not replace the
proof above or a Unit check.

\section{Relation to the local route}

The existing local route remains an alternative for the general case:
the field product theorem, local monogenic generation, finite
coefficient-and-denominator descent, monic rescaling, and finite
cover extraction. Lemma~\ref{lem:monicise} now shortens the construction
of the standard source; the direct lifting theorem supplies its
formally smooth target. Neither automatically supplies the remaining
global cover certificate.

For comparison, Stacks
\href{https://stacks.math.columbia.edu/tag/00UE}{Proposition 10.144.4}
gives the local standard-presentation theorem using a finite
subalgebra and a primitive element in a residue field. The point of
the present one-equation calculation is that this step can instead
be performed through explicit matrices and ideal certificates in
that case. Whether the same replacement can be made comparably
short for a general square Jacobian system remains open here.

\end{document}
